% Standalone research entry 215; batch 208--217.
% Original section material preserved from Pasted text(3).txt.
% Research additions: 27 September 2026.
% Compile twice with: lualatex 215.tex
% !TeX program = lualatex
% Compile twice: lualatex open_problems_500.tex
% All references and prose are embedded; no BibTeX database or external inputs.
\documentclass[11pt,a4paper]{article}
\usepackage[margin=24mm,headheight=14pt]{geometry}
\usepackage{fontspec}
\setmainfont{Latin Modern Roman}
\setsansfont{Latin Modern Sans}
\setmonofont{Latin Modern Mono}
\usepackage{amsmath,amssymb,mathtools}
\usepackage{unicode-math}
\setmathfont{Latin Modern Math}
\usepackage{microtype}
\usepackage{enumitem,xurl,needspace}
\usepackage[dvipsnames]{xcolor}
\usepackage{hyperref}
\hypersetup{unicode=true,colorlinks=true,linkcolor=MidnightBlue,urlcolor=MidnightBlue,
 pdftitle={500 Mathematical Problem Statements},pdfauthor={Source-annotated compilation},bookmarksnumbered=true}
\usepackage{bookmark}
\usepackage{tocloft}
\setlength{\cftsecnumwidth}{3em}
\cftsetpnumwidth{2.5em}
\cftsetrmarg{4em}
\usepackage{titlesec}
\titleformat{\section}{\Large\bfseries\raggedright}{\thesection}{0.7em}{}
\usepackage{fancyhdr}
\pagestyle{fancy}\fancyhf{}
\fancyhead[L]{\small\sffamily Mathematical problem statements}
\fancyhead[R]{\small Source edition 22}
\fancyfoot[C]{\thepage}
\setlength{\parindent}{0pt}
\setlength{\parskip}{5pt plus 1pt}
\setlength{\emergencystretch}{3em}
\setcounter{tocdepth}{1}
\setcounter{secnumdepth}{1}
\setlist[itemize]{leftmargin=3em,itemsep=5pt,topsep=4pt,parsep=0pt}
\allowdisplaybreaks
\newcommand{\N}{\mathbb{N}}
\newcommand{\Z}{\mathbb{Z}}
\newcommand{\Q}{\mathbb{Q}}
\newcommand{\R}{\mathbb{R}}
\newcommand{\C}{\mathbb{C}}
\newcommand{\F}{\mathbb{F}}
\newcommand{\A}{\mathbb{A}}
\newcommand{\PP}{\mathbb P}
\newcommand{\E}{\mathbb{E}}
\newcommand{\eps}{\varepsilon}
\newcommand{\dd}{\,\mathrm d}
\newcommand{\sref}[2]{\hyperlink{source:#1:#2}{[S#2]}}
\newcommand{\eref}[2]{\hyperlink{extra:#1:#2}{[E#2]}}
% Turn the next switch off for a shorter reading copy. The quotations preserve
% every exactTarget field, including qualifications omitted from a short summary.
\newif\ifshowcatalogue
\showcataloguetrue
\newcommand{\cataloguescope}[1]{\ifshowcatalogue
 \paragraph{Exact scope in the supplied catalogue.}
 \begin{quote}\small #1\end{quote}\fi}

% Standalone wrapper; no external inputs or bibliography files are required.
\fancyhead[L]{\small\sffamily Research notebook: section 215}
\fancyhead[R]{\small 27 September 2026}
\hypersetup{pdftitle={Research attempt 215},pdfauthor={Research notebook}}
\setlength{\parskip}{4pt plus 1pt}
\setlist[itemize]{before=\small\interlinepenalty10000,itemsep=3pt}
\begin{document}
\setcounter{section}{214}
% ================================================================
% problemId: problem.density-of-hyperbolicity-for-complex-polynomials
% source releaseRank: 215; source releaseStatus: open_with_solved_subcases
% exactTarget SHA-256: f383b93f602191b4e8268825addd163b157cbd8b4dc5f3617f943c18822204be
\clearpage
\section{\texorpdfstring{Density of Hyperbolicity for Complex Polynomials}{Density of Hyperbolicity for Complex Polynomials}}\label{215}
\subsection{Definitions and mathematical statement}
Let $\operatorname{Poly}_d=\{a_dz^d+\cdots+a_0:a_d\ne0\}$ with its coefficient topology. Regard a polynomial as a rational map of the Riemann sphere; infinity is its superattracting fixed point. A polynomial is hyperbolic when every critical orbit converges to an attracting periodic cycle, allowing escape to infinity in this sphere convention. The assertion is
\[
 \overline{\{f\in\operatorname{Poly}_d:f\text{ hyperbolic}\}}=\operatorname{Poly}_d
 \quad\text{for each }d\ge2.
\]
Equivalently, every coefficient neighborhood of every degree-$d$ polynomial contains such a map of the same degree. Requiring convergence only to finite attracting cycles would wrongly exclude escaping critical points and change the statement.
\par\smallskip\textit{Formulation and concept references:} \sref{215}{1}.
\cataloguescope{For each complex polynomial T of degree at least two, determine whether T can be approximated by polynomials of the same degree all of whose critical points converge under iteration to attracting periodic cycles, that is, whether hyperbolic polynomials are dense in the space of complex polynomials of each degree.}
\subsection{Short English statement}
Can every complex polynomial be approximated in its coefficients by one whose critical orbits all tend to attracting cycles, including the cycle at infinity?
% BEGIN RESEARCH ADDITION 215
\subsection{Research attempt}

\paragraph{Current frontier.}
The target is complex, fixed-degree density, including escape to infinity
\sref{215}{1}, Problem 11. The theorem of Kozlovski--Shen--van Strien
\eref{215}{1}, Theorem 1, concerns hyperbolicity of real polynomial dynamics
on the real line; its definition does not impose attraction of every
nonreal critical point. It must not be substituted for the present
complex statement. Dudko's December 2025 manuscript for ICM 2026 reports
bounded-type quadratic renormalization and local-connectivity results,
and explicitly isolates remaining unbounded satellite and interpolation
problems \eref{215}{2}, Sections 3.8--3.9 and 4.3--4.4.
These results suggest using quantitative critical relations, but do not
supply a perturbation theorem for all complex polynomials of every degree.
All estimates below are deductions made in this attempt; no global closing
lemma is being attributed to those sources.

\paragraph{Attack route.}
Preserve every critical orbit that already has a finite certificate of
attraction or escape, and try to close only the remaining critical orbits.
A simultaneous quantitative implicit-function argument could make those
critical points periodic and hence superattracting. The key issue is not
merely a return close to a critical point: the return error must be small
relative to the inverse parameter derivative and its variation. A second
route permits previously uncaptured critical points to escape. Two exact
quadratic examples below show why neither closing every critical point nor
assuming recurrence of every bad critical point is a valid universal plan.

\paragraph{Attempt.}
\textit{1. Robust finite capture certificates.}
For an attracting periodic point $z_*$ of period $q$, sufficiently small
closed disks $D$ about $z_*$ satisfy
\[
 f^q(D)\subset\operatorname{int}D,
 \qquad \sup_{z\in D}|(f^q)'(z)|<\theta<1.                 \tag{215.1}
\]
A critical orbit attracted to this cycle eventually enters the interior
of $D$. The finitely many iterates up to entry, the strict inclusion, and
the derivative inequality persist on a sufficiently small coefficient
neighborhood. The contraction theorem then supplies an attracting
periodic point for every perturbed map in that neighborhood. For escape,
a uniform escape radius $R$ exists on a small coefficient neighborhood
with the leading coefficient bounded away from zero. Entry into
$|z|>R$ is again a finite open condition. These arguments prove persistence
without taking an unjustified infinite intersection of neighborhoods.

It is enough for density to handle polynomials with simple critical
points: the discriminant of $f'$ is a nonzero polynomial in the
coefficients, so its nonvanishing locus is dense. First approximate a
multiple-critical-point polynomial within half the requested tolerance,
then solve the approximation problem for that simple-critical polynomial
within the remaining half. This reduction does \emph{not} make a further
parameter Jacobian automatically invertible.

\textit{2. A simultaneous closing lemma with explicit constants.}
Let $a\mapsto f_a$ be a holomorphic $k$-parameter family of polynomials of
the same degree, defined on a neighborhood of the closed ball
$\|a\|\le r$ in $\C^k$. Assume its remaining $k$ critical points have
holomorphic markings $c_i(a)$ there; all other critical points have the
persistent capture certificates above. For chosen positive integers $n_i$,
put
\[
 F_i(a)=f_a^{n_i}(c_i(a))-c_i(a),\qquad J=DF(0).
\]
Suppose $J$ is invertible, $L=\|J^{-1}\|$, and
\[
 \|J^{-1}F(0)\|\le r/2,
 \qquad \sup_{\|a\|\le r}\|DF(a)-J\|\le(2L)^{-1}.          \tag{215.2}
\]
Then a parameter $a_*$ with $\|a_*\|\le r$ satisfies $F(a_*)=0$.
Indeed $T(a)=a-J^{-1}F(a)$ has derivative norm at most $1/2$ throughout
the convex ball, and
\[
 \|T(a)\|\le\|J^{-1}F(0)\|+\tfrac12\|a\|\le r.
\]
Thus $T$ is a self-map contraction on a complete space. Its fixed point
also satisfies $\|a_*\|\le2\|J^{-1}F(0)\|$.

Every newly periodic critical point lies on a superattracting cycle,
because the multiplier contains a factor $f_{a_*}'(c_i(a_*))=0$.
The other critical points remain attracted or escaping. Consequently
$f_{a_*}$ is hyperbolic in precisely the sphere convention of the target.
This is a proved finite-parameter lemma. Applying it arbitrarily close to
every simple-critical polynomial requires the hypotheses (215.2) at
arbitrarily small parameter scales; that requirement remains unproved.
When there are fewer bad critical points than available parameters, a
$k$-dimensional parameter slice with invertible restricted Jacobian would
suffice, but its existence is also a hypothesis, not a genericity shortcut.

\textit{3. The derivative includes critical-point motion.}
The size and rank of $J$ can be tested rather than guessed. Write
$f_a=f+ah+O(a^2)$ along a parameter direction and let $c$ be a simple
critical point of $f$. Differentiating $f_a'(c(a))=0$ gives
\[
 Dc=-\frac{h'(c)}{f''(c)}.
\]
For $n\ge1$, the chain rule gives the exact identity
\[
 \begin{split}
 D\big(f_a^n(c(a))-c(a)\big)\big|_{a=0}
   ={}&\sum_{j=0}^{n-1}
      (f^{n-1-j})'(f^{j+1}(c))\,h(f^j(c))\\
     &+\frac{h'(c)}{f''(c)}.                              \end{split}\tag{215.3}
\]
The initial-point contribution to the derivative of the iterate is zero
because $(f^n)'(c)=0$; the last term comes from subtracting the moving
critical point. Omitting it can give the wrong rank or condition number.
For a concrete check take $f_b(z)=z^3+bz$, $3c^2+b=0$, $c\ne0$, and
$n=1$. Since $F=-2c^3-c$ and $dc/db=-1/(6c)$, direct differentiation gives
$dF/db=c+1/(6c)$, exactly as (215.3) with $h(z)=z$ predicts.
The reproduction script verifies this identity symbolically.

Long iterates create products of derivatives in (215.3). A tiny orbit-return
error can coexist with an almost singular $J$ or enormous variation of
$DF$ on the proposed parameter ball. Thus recurrence alone is not enough
to invoke a qualitative implicit-function theorem at a point where
$F(0)\ne0$. Conditions (215.2) retain precisely the quantitative control
that such an invocation would hide.

\textit{4. A false strengthening fails on an open hyperbolic set.}
For $f_c(z)=z^2+c$, consider
\[
 |c-1/10|<1/100,\qquad D=\{z:|z|\le1/5\}.
\]
For $z\in D$,
\[
 |f_c(z)|\le\frac1{25}+\frac{11}{100}=\frac3{20}<\frac15,
 \qquad |f_c'(z)|\le\frac25.
\]
The map is a contraction on $D$ and its only finite critical point $0$
is attracted to its fixed point there. These polynomials are hyperbolic.
If $0$ were periodic, contraction would force it to be that fixed point,
which would require $c=0$, excluded by the chosen parameter disk.
Hence even an open set of already hyperbolic maps need not contain a map
whose finite critical point is periodic. Density of critical-periodic
centers in the whole coefficient space is a false replacement for the
question. This is why the proposed strategy preserves captured critical
points instead of closing all of them.

\textit{5. A bad critical orbit need not return near its critical point.}
At $c=-2$, the critical orbit is $0,-2,2,2,\ldots$, so
$|f_c^n(0)-0|=2$ for all $n\ge1$. There is no small return residual
for the closing lemma at this parameter. Nevertheless, for real
$c=-a$ with $a>2$, one has
\[
 f_c^2(0)=a^2-a>a.
\]
For $x\ge a$, $x^2-a-x\ge a(a-2)>0$, so successive iterates grow
without bound. Thus $c<-2$ gives nearby hyperbolic maps by escape.
The example disproves an assumed universal recurrence input, not density.
A complete approach needs a rigorous alternative between controlled
critical closing, controlled escape, and other critical-relation
perturbations, including the neutral and infinitely renormalizable cases.

\paragraph{Stress test.}
The closing lemma is simultaneous: independently solving the $k$ scalar
equations and then combining their parameters is invalid. Its ball must
stay inside a common holomorphic marking domain and the capture
neighborhood for every already good critical point. No compactness
principle used here supplies those radii uniformly over long iterates.
The simple-critical reduction resolves multiplicities for a density
argument, but does not resolve singular Jacobians, nonrecurrence,
parabolic implosion, or recurrent critical orbits with badly controlled
scales. Polynomial perturbations also cannot independently alter one
orbit location by a localized smooth bump while keeping the degree fixed;
the holomorphic coefficient constraint is essential.

The role of renormalization is therefore specific: it would need to
supply quantitative parameter-scale control or a valid alternative to
(215.2), not merely assert that a dynamical return exists. The unfinished
regimes isolated in \eref{215}{2} are relevant warnings against assuming
such bounds. Hyperbolicity of a \emph{renormalization operator} in a
restricted class is not itself the requested density of hyperbolic
polynomials in all coefficient neighborhoods.

\paragraph{Outcome.}
\textbf{Outcome: Exploratory approach with unresolved gap.}

The attempt proves a robust-capture lemma and an explicit simultaneous
critical-closing criterion, with the full parameter derivative and two
exact tests rejecting stronger but false shortcuts. These are local tools,
not a solution of the density conjecture. The missing result is a
fixed-degree perturbation theorem guaranteeing controlled capture or
escape of every remaining critical orbit arbitrarily close to every
simple-critical polynomial. No argument here supplies that theorem, and
no counterexample to the original density statement has been constructed.
% END RESEARCH ADDITION 215
\subsection{Sources}
\begin{itemize}\raggedright
\item[\textnormal{[S1]}]\hypertarget{source:215:1}{} Steve Smale, Mathematical Problems for the Next Century, 1998.
\url{https://www.fim.uni-passau.de/fileadmin/dokumente/fakultaeten/fim/lehrstuhl/muller/SmaleProblems1998.pdf}.
{\small\textit{Catalogue formulation source.}}
% BEGIN RESEARCH REFERENCES 215
\item[\textnormal{[E1]}]\hypertarget{extra:215:1}{} O. Kozlovski, W. Shen and S. van Strien, ``Density of hyperbolicity in dimension one,'' \emph{Annals of Mathematics} 166 (2007), 145--182, Theorem 1 and its definition of real-line hyperbolicity.
\url{https://annals.math.princeton.edu/wp-content/uploads/annals-v166-n1-p04.pdf}.
\item[\textnormal{[E2]}]\hypertarget{extra:215:2}{} D. Dudko, ``On the MLC Conjecture and the Renormalization Theory in Complex Dynamics,'' arXiv:2512.24171 (30 December 2025), ICM 2026 manuscript, Sections 3.8--3.9 and 4.3--4.4, Theorems 3.3 and 4.1. The full MLC problem is not asserted solved there.
\url{https://arxiv.org/abs/2512.24171}.
% END RESEARCH REFERENCES 215
\end{itemize}

\end{document}
